\documentclass[11pt]{article}

\usepackage[margin=1in]{geometry}
\usepackage{amsmath,amssymb}
\usepackage{booktabs}
\usepackage{longtable}
\usepackage{array}
\usepackage{tabularx}
\usepackage{xcolor}
\usepackage{xurl}
\usepackage{hyperref}
\usepackage{enumitem}
\usepackage{microtype}
\usepackage{fancyhdr}

\hypersetup{
  colorlinks=true,
  linkcolor=blue!55!black,
  urlcolor=blue!65!black,
  citecolor=green!40!black,
  pdftitle={Fuel-Nozzle Technical Reference},
  pdfauthor={fuelnozzle project}
}

\pagestyle{fancy}
\fancyhf{}
\lhead{Fuel-Nozzle Technical Reference}
\rhead{Version 0.1.0}
\cfoot{\thepage}

\newcommand{\code}[1]{\nolinkurl{#1}}
\newcommand{\srcline}[2]{\nolinkurl{#1:#2}}
\newcommand{\unit}[1]{\,\mathrm{#1}}
\setlength{\emergencystretch}{3em}
\newcommand{\warningbox}[1]{%
  \begin{center}
  \fcolorbox{orange!60!black}{orange!8}{%
    \parbox{0.91\textwidth}{\textbf{Engineering caution.} #1}}
  \end{center}}

\title{Fuel-Nozzle Technical Reference\\
\large Jet-A Pressure-Swirl and Flashing LNG Models\\
\large Beginner-Oriented Equations, Algorithms, Source Locations, and Limitations}
\author{fuelnozzle project}
\date{Source map generated 2026-08-26}

\begin{document}
\maketitle

\begin{abstract}
This report explains the equations and software implementation of the
\code{fuelnozzle} package for a reader with no prior background in fluid mechanics,
thermodynamics, atomization, or Python. The package studies two fuel circuits:
a Jet-A pressure-swirl atomizer and a liquid-LNG nozzle that may remain liquid,
flash inside the nozzle, or flash after discharge. Every major equation is connected
to the current Python function and line where it is implemented. The report also
separates physical laws from numerical approximations and empirical calibration.
It should be read as documentation for a reduced-order screening tool, not as a
hardware certification method.
\end{abstract}

\tableofcontents
\newpage

\section{How to Read This Report}

A physical equation describes a relationship among measurable quantities. A
software implementation turns that relationship into arithmetic and decision logic.
The same equation can be implemented correctly while still being unsuitable for a
particular injector. Therefore, four questions must always be kept separate:

\begin{enumerate}[label=\arabic*.]
  \item \textbf{Was the equation entered correctly?} This is software verification.
  \item \textbf{Does the numerical method converge?} This is numerical verification.
  \item \textbf{Does the model match experiments?} This is validation.
  \item \textbf{Is the experiment similar enough to the intended LNG/Jet-A design?}
        This is the applicability question.
\end{enumerate}

Source references have the form \srcline{src/fuelnozzle/lng.py}{256}. They refer to
the source state on 2026-08-26. Function names are stable references; line numbers
will move when source files are edited.

\warningbox{The current code is a steady, zero/one-dimensional engineering screen.
It does not resolve the three-dimensional vortex, individual vapor bubbles, primary
breakup, combustion chemistry, flame stability, or structural integrity.}

\section{Basic Vocabulary for a New Reader}

\subsection{Pressure, temperature, and flow}

\begin{description}[style=nextline]
  \item[Pressure $P$ (Pa)] Force divided by area. One atmosphere is approximately
  $101325\unit{Pa}$. All code inputs are absolute pressures unless an external data
  source explicitly states otherwise.
  \item[Temperature $T$ (K)] Thermodynamic temperature. To convert Celsius to kelvin,
  add 273.15.
  \item[Mass flow $\dot m$ (kg/s)] Fuel mass passing a location each second.
  \item[Area $A$ ($\mathrm{m^2}$)] Flow-passage cross-sectional area.
  \item[Density $\rho$ ($\mathrm{kg/m^3}$)] Mass per unit volume.
  \item[Velocity $u$ (m/s)] Speed of the fluid.
  \item[Mass flux $G$ ($\mathrm{kg/(m^2\,s)}$)] Mass flow per area:
  \begin{equation}
    G = \frac{\dot m}{A} = \rho u.
    \label{eq:mass-flux-definition}
  \end{equation}
\end{description}

Mass flux is useful because a nozzle with twice the area carries twice the mass flow
when $G$ is unchanged.

\subsection{Liquid, vapor, saturation, and flashing}

A pure substance can exist as a liquid or vapor. At a given pressure below the
critical point, there is a saturation temperature. A liquid colder than saturation is
\emph{subcooled}. A liquid at saturation can begin boiling without further heating.

LNG is a mixture, so it does not have one boiling temperature. At a given pressure:

\begin{itemize}
  \item the \textbf{bubble temperature} is where the first vapor bubble can appear;
  \item the \textbf{dew temperature} is where the last liquid droplet disappears.
\end{itemize}

\emph{Flashing} is rapid vapor production caused by a pressure reduction. A liquid
can be thermodynamically able to flash but remain temporarily liquid because bubble
nucleation and growth require finite time. That temporary state is called
\emph{metastable}.

\subsection{Enthalpy, entropy, and quality}

\begin{description}[style=nextline]
  \item[Specific enthalpy $h$ (J/kg)] An energy-like state property useful for flow
  energy balances. In an adiabatic throttling process, $h$ is often approximately
  constant.
  \item[Specific entropy $s$ (J/(kg K))] A thermodynamic state property. An ideal
  adiabatic, reversible acceleration follows constant entropy.
  \item[Vapor quality $x$] Vapor mass divided by total liquid-plus-vapor mass. Thus
  $x=0$ is all liquid and $x=1$ is all vapor.
\end{description}

\subsection{Choking}

A flow is choked when reducing downstream pressure no longer increases mass flow.
For a flashing liquid, this may occur because vapor production makes the fluid highly
compressible. The current Tier 1 code calls a case choked when the maximum sampled
mass flux occurs at a pressure measurably above the outlet pressure.

\subsection{Atomization terms}

\begin{description}[style=nextline]
  \item[Pressure-swirl atomizer] A nozzle that sends liquid through tangential ports,
  forming a rotating liquid film and usually a hollow cone.
  \item[Air core] The low-density central region inside the rotating liquid annulus.
  \item[Film thickness] Radial thickness of liquid at the exit.
  \item[Cone angle] Angular width of the spray. This report distinguishes semi-angle
  from full cone angle. The code reports full cone angle.
  \item[SMD or $D_{32}$] Sauter mean diameter, the diameter of a hypothetical uniform
  spray with the same volume-to-surface-area ratio as the actual droplets.
  \item[$D_{10}$] Arithmetic mean droplet diameter. It is not interchangeable with SMD.
\end{description}

\section{Software Architecture and Data Flow}

For each flight-stage operating point, the software performs this sequence:

\begin{enumerate}
  \item Validate user inputs in \code{OperatingPoint}.
  \item Close the pump/feed/nozzle pressure budget.
  \item Create or reuse a CoolProp state for methane or an LNG composition.
  \item Optionally march the LNG through a heated feed line.
  \item Run LNG Tier 0, Tier 1, Tier 2, and Tier 3 in order.
  \item Independently run the Jet-A pressure-swirl calculation.
  \item Gather warnings instead of silently repairing inconsistent inputs.
  \item Integrate stage mass flows when all durations are available.
\end{enumerate}

The public study owner is \code{run_nozzle_study} at
\srcline{src/fuelnozzle/study.py}{75}. It creates one
\code{CoolPropLNGProvider} and reuses it for all operating points.

\section{Input Pressure Contract}

The user supplies compressor/combustor station values $P_3,T_3$, pump discharge
pressure, and nozzle pressure drop. The model currently treats $P_3$ as the receiving
pressure immediately downstream of each fuel nozzle.

The required nozzle inlet pressure is

\begin{equation}
  P_{n,\mathrm{in}} = P_3 + \Delta P_n.
  \label{eq:nozzle-inlet-pressure}
\end{equation}

The pressure remaining for pipes, bends, valves, and manifolds is

\begin{equation}
  \Delta P_{\mathrm{feed,available}}
  = P_{\mathrm{pump,out}} - P_{n,\mathrm{in}}.
  \label{eq:feed-pressure-budget}
\end{equation}

Equations~\eqref{eq:nozzle-inlet-pressure}--\eqref{eq:feed-pressure-budget} are
implemented by \code{fuel_pressure_budget} at
\srcline{src/fuelnozzle/operating.py}{55}, specifically lines 62--63.

If Equation~\eqref{eq:feed-pressure-budget} is negative, the point is hydraulically
infeasible. The code returns an error-level warning but does not alter $P_3$ or any
pressure input.

\warningbox{Real gas-turbine notation often uses $P_3$ for compressor discharge,
while pressure at the fuel nozzle can differ due to diffuser and liner losses. The
current interface does not separate these locations. Supply the pressure appropriate
to the nozzle receiving environment, or extend the data model.}

\section{Thermodynamic Property Model}

\subsection{LNG composition}

The user supplies nonnegative component mole fractions $z_i$. The input validator
normalizes them:

\begin{equation}
  z_i^{\mathrm{normalized}} = \frac{z_i}{\sum_j z_j}.
  \label{eq:composition-normalization}
\end{equation}

This occurs in \code{LNGComposition.validate_and_normalize} at
\srcline{src/fuelnozzle/models.py}{37}.

\subsection{CoolProp state objects}

\code{CoolPropLNGProvider.__init__} at
\srcline{src/fuelnozzle/properties.py}{29} creates one reusable CoolProp
\code{AbstractState}. Pure methane uses direct CoolProp input pairs. For a mixture,
CoolProp supports PT, PQ, and TQ state updates but not every PH or PS flash needed by
the nozzle path.

\subsection{Mixture PH and PS roots}

For a requested pressure $P$ and target enthalpy $h^*$, the code finds a temperature
that satisfies

\begin{equation}
  f_h(T) = h(P,T,\mathbf z)-h^* = 0.
  \label{eq:ph-root}
\end{equation}

For target entropy $s^*$, it solves

\begin{equation}
  f_s(T) = s(P,T,\mathbf z)-s^* = 0.
  \label{eq:ps-root}
\end{equation}

The common implementation is \code{_mixture_target_state} at
\srcline{src/fuelnozzle/properties.py}{152}. The residual is defined at line 160,
and SciPy's bracketed Brent method is called at line 210. The previous station's
temperature is used to find a nearby bracket; a wider scan is used only when local
continuation fails.

This approach is slower than a direct flash, but a bracketed root is robust and does
not require inventing mixture thermodynamics outside CoolProp.

\subsection{Molar quality to mass quality}

CoolProp's mixture quality is molar. The hydraulics require mass quality. Let $Q$ be
molar vapor fraction, $M_v$ vapor-phase molar mass, and $M_l$ liquid-phase molar
mass. Vapor and liquid masses for an arbitrary one-mole total basis are

\begin{align}
  m_v &= Q M_v, \\
  m_l &= (1-Q)M_l.
\end{align}

The mass quality is

\begin{equation}
  x = \frac{Q M_v}{Q M_v + (1-Q)M_l}.
  \label{eq:mass-quality-conversion}
\end{equation}

Implemented by \code{_mass_quality} at
\srcline{src/fuelnozzle/properties.py}{370}, with phase molar masses at lines
375--390 and Equation~\eqref{eq:mass-quality-conversion} at lines 391--393.

\subsection{Mixture transport fallback}

When CoolProp has no native LNG-mixture viscosity or conductivity, the provider uses
component values. Viscosity uses logarithmic mole mixing:

\begin{equation}
  \mu_{\mathrm{mix}} = \exp\left(\sum_i \hat z_i\ln\mu_i\right),
  \qquad \hat z_i=\frac{z_i}{\sum_{j\in\mathcal A}z_j},
  \label{eq:viscosity-mixing}
\end{equation}

where $\mathcal A$ is the set of components for which a value is available.
Conductivity and surface tension use linear mole mixing:

\begin{align}
  k_{\mathrm{mix}} &= \sum_i \hat z_i k_i, \\
  \sigma_{\mathrm{mix}} &= \sum_i \hat z_i \sigma_i.
  \label{eq:linear-transport-mixing}
\end{align}

Implemented by \code{_component_transport_fallback} at
\srcline{src/fuelnozzle/properties.py}{281}; Equations
\eqref{eq:viscosity-mixing}--\eqref{eq:linear-transport-mixing} are at lines
345--367.

If a pure component cannot be forced into a liquid PT state, its saturated-liquid
transport at the same temperature is attempted. Any omitted component is named in a
warning.

\warningbox{These transport mixing rules are engineering fallbacks, not a validated
LNG transport model. Density, enthalpy, entropy, and equilibrium still come from the
mixture EOS; only missing transport quantities are approximated.}

\section{Cryogenic LNG Feed Line}

\subsection{Segmentation}

The feed line is divided into equal axial segments. At each segment, the solver:

\begin{enumerate}
  \item reconstructs the state from local $P$ and $h$;
  \item estimates heat ingress;
  \item selects single- or two-phase pressure loss;
  \item subtracts pressure loss;
  \item adds heat to enthalpy.
\end{enumerate}

The owner is \code{solve_lng_feed_line} at
\srcline{src/fuelnozzle/feed.py}{63}.

\subsection{Pipe area and mass flux}

For inner diameter $D$,

\begin{equation}
  A_p = \frac{\pi D^2}{4}, \qquad G=\frac{\dot m}{A_p}.
  \label{eq:pipe-area-flux}
\end{equation}

Implemented at \srcline{src/fuelnozzle/feed.py}{94}--95.

\subsection{Segment energy update}

If heat leak per length is $q'$ and segment length is $\Delta z$, segment heat is
$q'\Delta z$. Dividing by mass flow gives the enthalpy rise:

\begin{equation}
  h_{j+1}=h_j+\frac{q'_j\Delta z}{\dot m}.
  \label{eq:feed-enthalpy-update}
\end{equation}

Implemented at \srcline{src/fuelnozzle/feed.py}{148}.

\subsection{Internal heat transfer}

The standard dimensionless groups are

\begin{align}
  \mathrm{Re} &= \frac{G D}{\mu}, \\
  \mathrm{Pr} &= \frac{c_p\mu}{k}, \\
  h_i &= \frac{\mathrm{Nu}\,k}{D}.
  \label{eq:feed-heat-groups}
\end{align}

ChEDL \code{ht} selects the internal-flow Nusselt correlation, and cylindrical
conduction combines internal convection, wall, insulation, and external convection.
The implementation is \code{_heat_leak_per_length} at
\srcline{src/fuelnozzle/feed.py}{230}, with the quantities at lines 244--267.
A measured heat leak in W/m bypasses this correlation chain.

\subsection{Single-phase pressure loss}

The dynamic pressure written in terms of mass flux is

\begin{equation}
  q_d=\frac{G^2}{2\rho}.
  \label{eq:dynamic-pressure}
\end{equation}

For Darcy friction factor $f$, distributed minor-loss coefficient $K$, segment
length $\Delta z$, and diameter $D$,

\begin{equation}
  \Delta P_{f+K}=\left(f\frac{\Delta z}{D}+K\right)\frac{G^2}{2\rho}.
  \label{eq:darcy-minor-loss}
\end{equation}

Static head is

\begin{equation}
  \Delta P_g=\rho g\Delta y.
  \label{eq:static-head}
\end{equation}

Implemented by \code{_liquid_segment_pressure_drop} at
\srcline{src/fuelnozzle/feed.py}{270}, with Equations
\eqref{eq:dynamic-pressure}--\eqref{eq:static-head} at lines 283--294.

\subsection{Two-phase pressure loss}

ChEDL \code{two_phase_dP} supplies frictional loss. The current selection is
\code{Kim_Mudawar} when $D\leq 6.22\unit{mm}$ and \code{Chisholm} for larger
lines. Static head uses homogeneous mixture density

\begin{equation}
  \rho_m=\left(\frac{1-x}{\rho_l}+\frac{x}{\rho_v}\right)^{-1}.
  \label{eq:homogeneous-density}
\end{equation}

Implemented by \code{_two_phase_segment_pressure_drop} at
\srcline{src/fuelnozzle/feed.py}{297}; correlation call and density occur at lines
315--333.

The current feed model does not add a separate two-phase acceleration pressure-drop
term. This is a known limitation when quality changes rapidly.

\section{LNG Tier 0: Is Flashing Thermodynamically Possible?}

\subsection{Subcooling margins}

The temperature subcooling margin is

\begin{equation}
  \Delta T_{sub}=T_{bubble}(P,\mathbf z)-T.
  \label{eq:temperature-margin}
\end{equation}

The pressure subcooling margin is

\begin{equation}
  \Delta P_{sub}=P-P_{bubble}(T,\mathbf z).
  \label{eq:pressure-margin}
\end{equation}

Positive values indicate equilibrium liquid margin. Implemented by
\code{screen_lng_flash} at \srcline{src/fuelnozzle/lng.py}{151}, specifically
lines 192--193.

\subsection{Equilibrium exit quality}

The nozzle inlet enthalpy is flashed isenthalpically to $P_3$. If the result is
two-phase, its mass quality is the equilibrium flash fraction. This is a
thermodynamic availability calculation, not a bubble-growth prediction.

\subsection{Equilibrium onset location}

The code assumes pressure changes linearly from nozzle inlet pressure $P_{in}$ to
$P_3$. If bubble pressure is $P_b$, the fractional axial position is

\begin{equation}
  \xi_{eq}=\frac{P_{in}-P_b}{P_{in}-P_3}.
  \label{eq:equilibrium-onset-location}
\end{equation}

Implemented at \srcline{src/fuelnozzle/lng.py}{223}. If $\xi_{eq}\geq0.98$, the
screen calls onset an exit event; otherwise it calls it internal.

\warningbox{Equation~\eqref{eq:equilibrium-onset-location} is a geometric screen.
The code does not solve the internal nozzle pressure profile from momentum and wall
friction.}

\section{LNG Tier 1: SPI and HEM Mass-Flow Bounds}

\subsection{Single-phase incompressible bound}

Bernoulli's equation gives ideal liquid speed $u=\sqrt{2\Delta P/\rho}$. Multiplying
by density gives

\begin{equation}
  G_{SPI}=\sqrt{2\rho_0\Delta P_n}.
  \label{eq:spi-flux}
\end{equation}

Implemented at \srcline{src/fuelnozzle/lng.py}{302}--306.

\subsection{Homogeneous equilibrium path}

HEM assumes liquid and vapor share velocity, pressure, and temperature and reach
local phase equilibrium instantly. The current path is isentropic:

\begin{equation}
  s(P)=s_0.
  \label{eq:isentropic-path}
\end{equation}

At each sampled pressure, stagnation enthalpy is converted to kinetic energy:

\begin{equation}
  u(P)=\sqrt{2\max\left(0,h_0-h(P,s_0)\right)}.
  \label{eq:hem-velocity}
\end{equation}

Mass flux follows:

\begin{equation}
  G_{HEM}(P)=\rho(P,s_0)u(P).
  \label{eq:hem-flux}
\end{equation}

Implemented by \code{solve_lng_equilibrium_flow} at
\srcline{src/fuelnozzle/lng.py}{256}; Equations
\eqref{eq:hem-velocity}--\eqref{eq:hem-flux} are lines 281--283.

\subsection{Critical condition and choking decision}

The sampled critical mass flux is

\begin{equation}
  G^*=\max_{P_3\leq P\leq P_{in}}G_{HEM}(P),
  \qquad P^*=\operatorname*{arg\,max}_P G_{HEM}(P).
  \label{eq:critical-flux}
\end{equation}

The software reports choking when

\begin{equation}
  P^*>P_3+\epsilon_P P_{in},
  \label{eq:choking-test}
\end{equation}

where default $\epsilon_P=0.005$. This logic is at lines 295--300 of
\code{solve_lng_equilibrium_flow}.

This is a numerical maximum criterion, not an explicit sonic-eigenvalue calculation.
Grid convergence is therefore a required verification test.

\subsection{Area and diameter sizing}

For discharge coefficient $C_d$ and selected operating mass flux $G_{op}$,

\begin{equation}
  A_{req}=\frac{\dot m}{C_dG_{op}}.
  \label{eq:lng-area}
\end{equation}

For $N$ equal circular holes, diameter per hole is

\begin{equation}
  D_{req}=\sqrt{\frac{4A_{req}}{\pi N}}.
  \label{eq:lng-diameter}
\end{equation}

Implemented at \srcline{src/fuelnozzle/lng.py}{310}--315.

The discharge coefficient is user-specified (default 0.80). HEM does not itself
predict vena-contracta and wall-loss effects.

\section{LNG Tier 2: Finite-Rate Relaxation Screen}

\subsection{Why a relaxation model is needed}

A short nozzle may not allow enough time for equilibrium vapor to form. The conceptual
continuous model is

\begin{equation}
  \frac{Dx}{Dt}=\frac{x_{eq}-x}{\tau},
  \label{eq:continuous-relaxation}
\end{equation}

where $\tau$ is a fitted relaxation time. The implementation uses the exact update for
constant $x_{eq}$ over one time step:

\begin{equation}
  x_{j+1}=x_j+\left(1-e^{-\Delta t_j/\tau}\right)
  \left(x_{eq,j+1}-x_j\right).
  \label{eq:discrete-relaxation}
\end{equation}

\subsection{Nucleation delay}

Relaxation is allowed only below

\begin{equation}
  P_{nuc}=\max\left(0,P_{bubble}-\Delta P_{delay}\right).
  \label{eq:nucleation-threshold}
\end{equation}

A larger $\Delta P_{delay}$ lets liquid remain metastable to lower pressure.
Equation~\eqref{eq:nucleation-threshold} is at
\srcline{src/fuelnozzle/lng.py}{385}--389.

\subsection{Residence time}

Pressure samples are mapped to equally spaced axial positions. For segment length
$\Delta z$ and average endpoint speed,

\begin{equation}
  \overline u_j=\max\left(0.5,\frac{u_j+u_{j+1}}{2}\right),
  \qquad \Delta t_j=\frac{\Delta z}{\overline u_j}.
  \label{eq:residence-step}
\end{equation}

The floor of $0.5\unit{m/s}$ prevents division by zero at the inlet. Implemented at
\srcline{src/fuelnozzle/lng.py}{405}--408. Equation
\eqref{eq:discrete-relaxation} is at lines 410--416.

\subsection{Actual mixture density}

For $x>0$, the code uses the homogeneous density from
Equation~\eqref{eq:homogeneous-density}. For $x=0$, it retains compressed-liquid inlet
density. Implemented by \code{_relaxed_mixture_density} at
\srcline{src/fuelnozzle/lng.py}{539}--552.

\subsection{Bounded mass flux}

The direct estimate is $G_{raw}=\rho_{actual}u_{HEM}$. The implementation prevents
this reduced model from leaving the HEM/SPI bracket:

\begin{equation}
  G_{Tier2}=\min\left[G_{SPI},\max\left(G_{HEM},G_{raw}\right)\right].
  \label{eq:tier2-bounded-flux}
\end{equation}

Implemented at \srcline{src/fuelnozzle/lng.py}{428}--446.

\warningbox{Tier 2 does not solve coupled mass, momentum, energy, and phase-relaxation
equations. It relaxes quality along the Tier 1 pressure/velocity path and applies a
physical bracket. It should be called a finite-rate screening model, not a validated
full HRM solver.}

\section{LNG Tier 3: Spray Regime and Calibration}

\subsection{Pressure ratio, superheat, and Jakob number}

The saturation-to-chamber pressure ratio is

\begin{equation}
  R_p=\frac{P_{bubble}(T_{in},\mathbf z)}{P_3}.
  \label{eq:pressure-ratio}
\end{equation}

Superheat relative to the bubble temperature at $P_3$ is

\begin{equation}
  \Delta T_{sh}=\max\left(0,T_{in}-T_{bubble}(P_3,\mathbf z)\right).
  \label{eq:superheat}
\end{equation}

Latent heat estimate and Jakob number are

\begin{align}
  h_{fg} &= h_{dew}(P_3)-h_{bubble}(P_3), \\
  \mathrm{Ja} &= \frac{c_p\Delta T_{sh}}{h_{fg}}.
  \label{eq:jakob}
\end{align}

Implemented by \code{solve_lng_flash_spray} at
\srcline{src/fuelnozzle/spray.py}{102}; equations occur at lines 126--143.

\subsection{Regime classifier}

\code{_classify_regime} begins at \srcline{src/fuelnozzle/spray.py}{225}.
The classifier first honors upstream, no-flash, and external-flash results from lower
tiers. Remaining cases are fully flashing when any of these default conditions is true:

\begin{equation}
  x_{exit}\geq0.10 \quad\text{or}\quad R_p\geq3.0
  \quad\text{or}\quad \mathrm{Ja}\geq0.10.
  \label{eq:full-flash-threshold}
\end{equation}

Otherwise the case is transitional. These thresholds are exposed in
\code{FlashSpraySettings} at \srcline{src/fuelnozzle/spray.py}{35}; they are not
universal physical boundaries.

\subsection{Calibrated LNG SMD}

No SMD is produced without calibration. With a reference SMD, pressure ratio, and
flash quality, the code uses

\begin{equation}
  D_{32}=D_{32,ref}
  \left(\frac{R_p}{R_{p,ref}}\right)^{-a}
  \left(\frac{x_{drive}}{x_{ref}}\right)^{-b},
  \label{eq:lng-smd-calibration}
\end{equation}

where

\begin{equation}
  x_{drive}=\max\left(10^{-6},x_{exit},x_{eq,P3}\right).
  \label{eq:driving-quality}
\end{equation}

The uncertainty interval is

\begin{equation}
  [D_{32}(1-U_D),\;D_{32}(1+U_D)].
  \label{eq:lng-smd-range}
\end{equation}

Implemented at \srcline{src/fuelnozzle/spray.py}{168}--181.

\subsection{Calibrated cone angle}

The full cone-angle estimate is

\begin{equation}
  \Theta=\operatorname{clip}_{[0,180]}
  \left[\Theta_{ref}+K_\Theta\ln\left(\frac{R_p}{R_{p,ref}}\right)\right].
  \label{eq:lng-cone-calibration}
\end{equation}

The reported range adds and subtracts a calibration uncertainty and clips it to
$[0,180]$ degrees. Implemented at
\srcline{src/fuelnozzle/spray.py}{182}--189.

\warningbox{Equations~\eqref{eq:lng-smd-calibration} and
\eqref{eq:lng-cone-calibration} are calibration forms selected for transparent
screening. They are not universal flash-atomization correlations.}

\section{Jet-A Pressure-Swirl Model}

\subsection{Jet-A properties}

Jet-A is a variable mixture rather than one pure chemical. The code therefore accepts
density, viscosity, surface tension, and optional vapor pressure from a measured or
declared source. \code{JetAPropertyTable.at_temperature} linearly interpolates a
monotonic temperature table at \srcline{src/fuelnozzle/jet_a.py}{64}. Outside the
table it uses an endpoint and emits a warning.

\subsection{Ideal speed and exit sizing}

The Bernoulli ideal speed is

\begin{equation}
  u_{ideal}=\sqrt{\frac{2\Delta P_n}{\rho}}.
  \label{eq:jet-ideal-velocity}
\end{equation}

The orifice relation is

\begin{equation}
  \dot m=C_d\rho A_o u_{ideal}
  =C_d A_o\sqrt{2\rho\Delta P_n}.
  \label{eq:jet-orifice-flow}
\end{equation}

Solving for area and diameter gives

\begin{align}
  A_{req}&=\frac{\dot m}{C_d\rho u_{ideal}}, \\
  D_{req}&=\sqrt{\frac{4A_{req}}{\pi}}.
  \label{eq:jet-sizing}
\end{align}

Implemented by \code{solve_jet_a_pressure_swirl} at
\srcline{src/fuelnozzle/jet_a.py}{143}; equations occur at lines 168--180.

\subsection{Tangential port velocity and angular momentum}

For $N_p$ circular ports of diameter $D_p$,

\begin{equation}
  A_p=N_p\frac{\pi D_p^2}{4}, \qquad
  u_p=\frac{\dot m}{\rho A_p}.
  \label{eq:port-flow}
\end{equation}

A total exit speed is formed using velocity coefficient $C_v$:

\begin{equation}
  u_{tot}=C_v u_{ideal}.
  \label{eq:total-exit-speed}
\end{equation}

The unconstrained tangential speed follows a simple angular-momentum transfer:

\begin{equation}
  u_{\theta,raw}=\eta_L u_p\frac{r_t}{r_o},
  \label{eq:angular-momentum-closure}
\end{equation}

where $r_t$ is port tangency radius, $r_o$ is exit radius, and $\eta_L$ is an
empirical efficiency. Implemented at \srcline{src/fuelnozzle/jet_a.py}{183}--193.

The code caps tangential speed using total pressure energy and a minimum axial speed
needed for continuity. Axial speed is then

\begin{equation}
  u_z=\sqrt{\max\left(10^{-12},u_{tot}^2-u_\theta^2\right)}.
  \label{eq:axial-velocity}
\end{equation}

Implemented at \srcline{src/fuelnozzle/jet_a.py}{194}--217.

\subsection{Air core and film thickness}

The liquid annular area required by continuity is

\begin{equation}
  A_l=\frac{\dot m}{\rho u_z}.
  \label{eq:liquid-annulus-area}
\end{equation}

For exit radius $r_o$, the air-core radius is

\begin{equation}
  r_a=\sqrt{r_o^2-\frac{A_l}{\pi}},
  \label{eq:air-core-radius}
\end{equation}

and film thickness is

\begin{equation}
  t_f=r_o-r_a.
  \label{eq:film-thickness}
\end{equation}

Implemented at \srcline{src/fuelnozzle/jet_a.py}{218}--234. If $A_l\geq A_o$,
the code sets no air core and returns an error warning.

\subsection{Cone angle}

The nominal full cone angle is twice the velocity-vector semi-angle:

\begin{equation}
  \Theta=2\tan^{-1}\left(\frac{u_\theta}{u_z}\right).
  \label{eq:jet-cone-angle}
\end{equation}

Implemented at \srcline{src/fuelnozzle/jet_a.py}{236}. This is simpler than a
complete Abramovich/Giffen-Muraszewski internal-flow solution.

\subsection{Cavitation screen}

Centrifugal velocity is assumed to depress internal pressure by dynamic pressure:

\begin{equation}
  P_{min}=P_{n,in}-\frac{1}{2}\rho u_\theta^2.
  \label{eq:jet-minimum-pressure}
\end{equation}

The vapor-pressure margin is

\begin{equation}
  \Delta P_{cav}=P_{min}-P_{vap}(T).
  \label{eq:cavitation-margin}
\end{equation}

Implemented at \srcline{src/fuelnozzle/jet_a.py}{237}--244. Negative margin
triggers an error-level warning. This is a screen, not a resolved cavitation model.

\subsection{Dimensionless atomization groups}

Using film thickness as length scale,

\begin{align}
  \mathrm{Re}_f &= \frac{\rho u_{tot}t_f}{\mu}, \\
  \mathrm{We}_f &= \frac{\rho u_{tot}^2t_f}{\sigma}, \\
  \mathrm{Oh}_f &= \frac{\mu}{\sqrt{\rho\sigma t_f}}.
  \label{eq:jet-dimensionless}
\end{align}

Implemented at \srcline{src/fuelnozzle/jet_a.py}{262}--275.
Reynolds number compares inertia to viscosity, Weber number compares inertia to
surface tension, and Ohnesorge number compares viscosity to inertial-capillary forces.

Ambient air density is estimated as an ideal gas:

\begin{equation}
  \rho_a=\frac{P_3}{R_aT_3}, \qquad R_a=287.05\unit{J/(kg\,K)}.
  \label{eq:ambient-air-density}
\end{equation}

Implemented at \srcline{src/fuelnozzle/jet_a.py}{276}--278.

\subsection{Calibrated Jet-A SMD}

The code defines a capillary sheet scale

\begin{equation}
  \ell_c=\sqrt{\frac{\sigma t_f}{\rho u_{tot}^2}}
  \label{eq:capillary-sheet-scale}
\end{equation}

and calibrated SMD

\begin{equation}
  D_{32}=C_{SMD}\ell_c(1+3\mathrm{Oh}_f).
  \label{eq:jet-smd-calibration}
\end{equation}

Implemented at \srcline{src/fuelnozzle/jet_a.py}{294}--303. No SMD is returned
unless $C_{SMD}$ is supplied.

\warningbox{The current geometry fields for swirl-chamber radius, swirl-chamber
length, and exit-orifice length are validated inputs but do not enter the hydraulic
equations. A future named pressure-swirl model must use them or remove them from the
claimed physics.}

\section{Mission-Level Integration}

If every operating point has duration $\Delta t_i$ and multiplier $N_i$, integrated
fuel mass is

\begin{equation}
  m_f=\sum_i \dot m_{f,i}\Delta t_i N_i.
  \label{eq:mission-mass}
\end{equation}

Implemented by \code{_integrate_fuel_masses} at
\srcline{src/fuelnozzle/study.py}{170}; Jet-A and LNG sums are lines 175--184.

For external mission total $m_{ref}$, mismatch is

\begin{equation}
  e_{rel}=\frac{|m_f-m_{ref}|}{\max(m_{ref},10^{-12})}.
  \label{eq:mission-relative-error}
\end{equation}

Implemented by \code{_mission_mass_warnings} at
\srcline{src/fuelnozzle/study.py}{188}, equation at line 199. The default warning
tolerance is 5\%.

\section{Complete Function-to-Source Map}

\footnotesize
\begin{longtable}{>{\raggedright\arraybackslash}p{0.31\textwidth}
                  >{\raggedright\arraybackslash}p{0.22\textwidth}
                  >{\raggedright\arraybackslash}p{0.39\textwidth}}
\toprule
Function or class & Source start & Responsibility \\
\midrule
\endfirsthead
\toprule
Function or class & Source start & Responsibility \\
\midrule
\endhead
\code{OperatingPoint} & \srcline{src/fuelnozzle/operating.py}{12} & Per-stage user inputs \\
\code{fuel_pressure_budget} & \srcline{src/fuelnozzle/operating.py}{55} & Pump/feed/nozzle pressure closure \\
\code{LNGComposition} & \srcline{src/fuelnozzle/models.py}{28} & Normalize LNG composition \\
\code{CoolPropLNGProvider.__init__} & \srcline{src/fuelnozzle/properties.py}{29} & Reusable EOS state creation \\
\code{state_pt} & \srcline{src/fuelnozzle/properties.py}{60} & Pressure-temperature flash \\
\code{state_ph} & \srcline{src/fuelnozzle/properties.py}{72} & Pressure-enthalpy flash/root \\
\code{state_ps} & \srcline{src/fuelnozzle/properties.py}{96} & Pressure-entropy flash/root \\
\code{state_pq} & \srcline{src/fuelnozzle/properties.py}{120} & Pressure-quality saturation state \\
\code{state_tq} & \srcline{src/fuelnozzle/properties.py}{132} & Temperature-quality saturation state \\
\code{_mixture_target_state} & \srcline{src/fuelnozzle/properties.py}{152} & Bracketed mixture PH/PS root \\
\code{_component_transport_fallback} & \srcline{src/fuelnozzle/properties.py}{281} & Approximate missing mixture transport \\
\code{_mass_quality} & \srcline{src/fuelnozzle/properties.py}{370} & Convert molar to mass quality \\
\code{solve_lng_feed_line} & \srcline{src/fuelnozzle/feed.py}{63} & Segmented cryogenic feed solver \\
\code{_heat_leak_per_length} & \srcline{src/fuelnozzle/feed.py}{230} & Measured or calculated heat ingress \\
\code{_liquid_segment_pressure_drop} & \srcline{src/fuelnozzle/feed.py}{270} & Single-phase friction/minor/static loss \\
\code{_two_phase_segment_pressure_drop} & \srcline{src/fuelnozzle/feed.py}{297} & ChEDL two-phase friction and static loss \\
\code{screen_lng_flash} & \srcline{src/fuelnozzle/lng.py}{151} & Tier 0 equilibrium phase screen \\
\code{solve_lng_equilibrium_flow} & \srcline{src/fuelnozzle/lng.py}{256} & Tier 1 SPI/HEM sizing and choking \\
\code{solve_lng_relaxation_flow} & \srcline{src/fuelnozzle/lng.py}{363} & Tier 2 finite-rate flashing screen \\
\code{_relaxed_mixture_density} & \srcline{src/fuelnozzle/lng.py}{539} & Homogeneous actual-quality density \\
\code{solve_lng_flash_spray} & \srcline{src/fuelnozzle/spray.py}{102} & Tier 3 spray outputs and calibration \\
\code{_classify_regime} & \srcline{src/fuelnozzle/spray.py}{225} & Rule-based flash-spray class \\
\code{JetAPropertyTable.at_temperature} & \srcline{src/fuelnozzle/jet_a.py}{64} & Jet-A property interpolation \\
\code{solve_jet_a_pressure_swirl} & \srcline{src/fuelnozzle/jet_a.py}{143} & Jet-A hydraulic and spray screen \\
\code{run_nozzle_study} & \srcline{src/fuelnozzle/study.py}{75} & Whole-envelope orchestration \\
\code{_integrate_fuel_masses} & \srcline{src/fuelnozzle/study.py}{170} & Stage-to-mission fuel integration \\
\code{_mission_mass_warnings} & \srcline{src/fuelnozzle/study.py}{188} & External mission consistency check \\
\bottomrule
\end{longtable}
\normalsize

\section{Existing Verification Coverage}

The current suite has 21 tests. The main test starts are:

\begin{itemize}
  \item properties: \srcline{tests/test_properties.py}{7}, 19, and 51;
  \item feed line: \srcline{tests/test_feed.py}{31} and 48;
  \item LNG Tier 0/1: \srcline{tests/test_lng_equilibrium.py}{37}, 50, 63, and 81;
  \item LNG Tier 2: \srcline{tests/test_lng_relaxation.py}{36}, 52, and 72;
  \item LNG Tier 3: \srcline{tests/test_spray.py}{35}, 54, and 80;
  \item Jet-A: \srcline{tests/test_jet_a.py}{41}, 67, and 83;
  \item full study: \srcline{tests/test_study.py}{47}, 70, and 86.
\end{itemize}

These tests confirm code paths, round trips, bounds, warning behavior, and basic limiting
cases. Required next verification includes pressure-grid convergence, Tier 2 axial-grid
convergence, independent analytical values, energy/momentum residuals, and property tables.

\section{Validation Evidence and Planned Use}

\subsection{Critical flow}

Hammer et al.~\cite{hammer2022} provide modern CO2 orifice/nozzle choking data.
CO2 is not an LNG surrogate; it tests the HEM and non-equilibrium mathematical
structure for a strongly real fluid.

De Lorenzo et al.~\cite{delorenzo2017} benchmark delayed-equilibrium and classic
critical-flow models against multiple water experiments, including Super Moby Dick.
This is the primary source for testing relaxation-model behavior.

Kim and O'Neal~\cite{kimoneal1995} and the Payne/O'Neal compilation provide
short-orifice refrigerant evidence for metastability and geometry dependence.

\subsection{Cryogenic spray}

DaRUS-2527~\cite{darus2527} contains 22 LN2 shadowgraph cases and condition data.
It is the preferred Tier 3 regime/angle dataset. Rees et al.~\cite{reesmorphology}
provide interpretation.

DaRUS-2076~\cite{darus2076} and Rees et al.~\cite{rees2020} provide raw local
$D_{10}$ and velocity data. These validate cryogenic spray fields but do not directly
validate LNG $D_{32}$.

Gaertner et al.~\cite{gaertner2020} provide a numerical/experimental cryogenic
comparison. The correct article number is 103360.

\subsection{Pressure swirl and Jet-A}

Lacava et al.~\cite{lacava2004} provide a complete, open, four-port water atomizer
benchmark. At the 4-atm design point the reported experimental values are
$\dot m=6.13\unit{g/s}$, $C_d=0.2728$, semi-angle $34.5^\circ$, and
$D_{32}=80.83\unit{\mu m}$. Water data verify model form, not Jet-A transfer.

Dodge and Biaglow~\cite{dodge1986} provide direct Jet-A SMD and distribution trends
with elevated air pressure and temperature. The NASA record is metadata-only, so a
full paper must be transcribed before quantitative tests are created.

\section{Known Deviations and Technical Debt}

\begin{enumerate}
  \item $P_3$ is treated as nozzle back pressure; no separate liner pressure loss exists.
  \item The feed model is steady and uniform. It omits pump dynamics, valve transients,
  conjugate transient wall storage, and a separate two-phase acceleration loss.
  \item Tier 0 onset position assumes a linear pressure profile.
  \item Tier 1 samples an isentropic thermodynamic path; it does not solve an
  area/friction-coupled nozzle differential equation.
  \item Choking is a sampled mass-flux maximum rather than an explicit sonic condition.
  \item Tier 2 is a quality-relaxation screen on the Tier 1 path, not a full HRM PDE/ODE
  conservation solver.
  \item Tier 2 mass flux is clamped to HEM/SPI bounds. This is a guardrail, not a derived
  non-equilibrium critical-flow law.
  \item Tier 3 regime thresholds and calibrated SMD/cone equations require validation.
  \item LNG transport mixing is approximate where CoolProp has no native mixture model.
  \item Jet-A swirl chamber radius, chamber length, and exit length are not yet used by
  the hydraulic closure.
  \item The Jet-A cone angle uses a velocity-vector approximation rather than a complete
  pressure-swirl theory.
  \item No automated parameter uncertainty propagation is implemented.
  \item No direct adapter to the separate mission repository is implemented; only generic
  operating points and mission-total checks are supported.
  \item CFD export is an in-memory boundary record, not a solver-specific file writer.
\end{enumerate}

\section{Recommended Next Steps}

\begin{enumerate}
  \item Add numerical-convergence tests before fitting any empirical parameter.
  \item Create immutable raw-data and derived-data folders with source checksums.
  \item Transcribe one Hammer CO2 orifice series and all Lacava curves first.
  \item Generalize the property protocol so the same Tier 1 kernel can run CO2, R134a,
  water, nitrogen, methane, and LNG without changing equations.
  \item Fit Tier 2 on declared training cases and evaluate held-out geometries/states.
  \item Load DaRUS datasets directly rather than digitizing available CSV/XLSX values.
  \item Add named pressure-swirl correlation adapters and use currently unused geometry.
  \item Add direct LCH4/LNG validation only when composition, inlet state, geometry,
  back pressure, mass flow, and uncertainty are all available.
\end{enumerate}

\section{Building This Report}

The project uses Tectonic through Pixi:

\begin{verbatim}
pixi install
pixi run doc
\end{verbatim}

The task is defined in \code{pixi.toml}. Tectonic downloads required TeX packages and
writes the PDF next to this source file.

\begin{thebibliography}{99}

\bibitem{setzmann1991}
U. Setzmann and W. Wagner,
``A New Equation of State and Tables of Thermodynamic Properties for Methane
Covering the Range from the Melting Line to 625 K at Pressures up to 1000 MPa,''
\emph{Journal of Physical and Chemical Reference Data}, vol. 20, no. 6,
pp. 1061--1155, 1991.
\url{https://doi.org/10.1063/1.555898}.

\bibitem{bell2014}
I. H. Bell et al.,
``Pure and Pseudo-pure Fluid Thermophysical Property Evaluation and the
Open-Source Thermophysical Property Library CoolProp,''
\emph{Industrial \& Engineering Chemistry Research}, vol. 53, 2014.
\url{https://doi.org/10.1021/ie4033999}.

\bibitem{hammer2022}
M. Hammer, H. Deng, A. Austegard, A. M. Log, and S. T. Munkejord,
``Experiments and modelling of choked flow of CO2 in orifices and nozzles,''
\emph{International Journal of Multiphase Flow}, vol. 156, article 104201, 2022.
\url{https://doi.org/10.1016/j.ijmultiphaseflow.2022.104201}.

\bibitem{delorenzo2017}
M. De Lorenzo, Ph. Lafon, J.-M. Seynhaeve, and Y. Bartosiewicz,
``Benchmark of Delayed Equilibrium Model (DEM) and classic two-phase critical
flow models against experimental data,''
\emph{International Journal of Multiphase Flow}, vol. 92, pp. 112--130, 2017.
\url{https://doi.org/10.1016/j.ijmultiphaseflow.2017.03.004}.

\bibitem{kimoneal1995}
Y. Kim and D. L. O'Neal,
``A comparison of critical flow models for estimating two-phase flow of HCFC22
and HFC134a through short-tube orifices,''
\emph{International Journal of Refrigeration}, 1995.
\url{https://doi.org/10.1016/0140-7007(95)93785-I}.

\bibitem{darus2527}
A. Rees and M. Oschwald,
``Evolution of flash boiling liquid nitrogen sprays by means of high-speed
shadowgraphy,'' DaRUS, V1, 2022.
\url{https://doi.org/10.18419/DARUS-2527}.

\bibitem{reesmorphology}
A. Rees, H. Salzmann, J. Sender, and M. Oschwald,
``About the Morphology of Flash Boiling Liquid Nitrogen Sprays,''
\emph{Atomization and Sprays}, vol. 30, no. 10, pp. 713--740, 2020.
\url{https://doi.org/10.1615/AtomizSpr.2020035265}.

\bibitem{darus2076}
A. Rees and M. Oschwald,
``Experimental data of flash boiling liquid nitrogen sprays (Test Case IN-1),''
DaRUS, V1, 2021.
\url{https://doi.org/10.18419/DARUS-2076}.

\bibitem{rees2020}
A. Rees, L. Araneo, H. Salzmann, G. Lamanna, J. Sender, and M. Oschwald,
``Droplet velocity and diameter distributions in flash boiling liquid nitrogen
jets by means of phase Doppler diagnostics,''
\emph{Experiments in Fluids}, vol. 61, article 182, 2020.
\url{https://doi.org/10.1007/s00348-020-03020-7}.

\bibitem{gaertner2020}
J. W. Gaertner, A. Kronenburg, A. Rees, J. Sender, M. Oschwald, and G. Lamanna,
``Numerical and experimental analysis of flashing cryogenic nitrogen,''
\emph{International Journal of Multiphase Flow}, vol. 130, article 103360, 2020.
\url{https://doi.org/10.1016/j.ijmultiphaseflow.2020.103360}.

\bibitem{lacava2004}
P. T. Lacava, D. Bastos-Netto, and A. P. Pimenta,
``Design Procedure and Experimental Evaluation of Pressure-Swirl Atomizers,''
24th International Congress of the Aeronautical Sciences, 2004.
\url{https://www.icas.org/icas_archive/ICAS2004/PAPERS/097.PDF}.

\bibitem{dodge1986}
L. G. Dodge and J. A. Biaglow,
``Effect of elevated temperature and pressure on sprays from simplex swirl
atomizers,'' \emph{Journal of Engineering for Gas Turbines and Power},
vol. 108, ASME paper 85-GT-58, 1986.
\url{https://ntrs.nasa.gov/citations/19860037997}.

\end{thebibliography}

\end{document}
